\documentclass[11pt]{article}
\usepackage[margin=1in]{geometry}
\usepackage{amsmath}
\usepackage{natbib}
\usepackage{hyperref}

\title{Force-dependent reaction pathways\\in a model system}
\author{Demo Manuscript}
\date{}

\begin{document}
\maketitle

\begin{abstract}
Mechanical force can alter the rates of chemical reactions.
In this study, we investigate how a model reaction responds to force.
We find that the reaction gets faster when more force is applied.
The results are consistent with a force-dependent activation barrier.
\end{abstract}

\section{Introduction}
The relationship between molecular structure and mechanical response
motivates the study of force-dependent reactivity. A simple model
provides a useful starting point for comparing reaction pathways.
This document contains synthetic examples for testing a writing interface;
the values and conclusions below are not experimental findings.
Further details are provided in the demonstration notes~\citep{demo}.

\section{Model and methods}
We consider a one-dimensional reaction coordinate and represent the
effect of an applied force using a linear coupling term. The
force-dependent activation free energy is written as
\begin{equation}
  \Delta G^{\ddagger}(F) = \Delta G^{\ddagger}(0) - F\Delta x^{\ddagger},
  \label{eq:barrier}
\end{equation}
where $F$ is the applied force and $\Delta x^{\ddagger}$ is the distance
between the initial state and the transition state along the force axis.

The model uses an illustrative temperature of $T = 298\,\mathrm{K}$ and
an illustrative transition distance of $\Delta x^{\ddagger} = 0.12\,\mathrm{nm}$.
We examine forces from $0$ to $200\,\mathrm{pN}$ in increments of
$25\,\mathrm{pN}$. These values are used only to exercise the preview.

\section{Results and discussion}
Equation~\ref{eq:barrier} predicts a lower activation barrier as the applied
force increases for a positive transition distance.
This is due to the fact that the force makes the barrier lower.
The corresponding rate depends exponentially on the activation free energy:
\begin{equation}
  k(F) = k(0)\exp\!\left(\frac{F\Delta x^{\ddagger}}{k_{\mathrm{B}}T}\right).
  \label{eq:rate}
\end{equation}
This expression assumes that the transition distance remains constant
over the force range of interest. Changes in mechanism or in the
underlying free-energy surface would require a more detailed model.

\section{Conclusion}
We have shown that force controls the reaction mechanism.
Future work will compare the simple model with a more complete
description of the free-energy surface.

\bibliographystyle{plainnat}
\bibliography{references}
\end{document}
